\section{Conclusion}
% This work revisits CPU-based acceleration of AlphaFold2 in light of recent architectural innovations in many-core processors. By leveraging the 72F’s matrix units and integrating quantization into the computation pipeline, MatrixFold achieves competitive performance with GPU implementations and, crucially, outperforms the A800 GPU on long-sequence inference where memory capacity becomes a limiting factor for GPUs. Our NUMA-aware communication and adaptive parallelism strategies ensure efficient scaling across diverse sequence lengths. Experimental results on CASP14 confirm that both FP16 and INT8 quantization maintain reliable accuracy while substantially reducing inference latency. Beyond AlphaFold2, the techniques proposed in MatrixFold are applicable to other Evoformer-based models such as RoseTTAFold and ESMFold, highlighting the broader impact of CPU-based quantized inference. This study demonstrates that CPUs with matrix extensions are not only viable but can exceed top-tier GPUs in large-scale protein structure prediction workloads.
In this paper we demonstrate that through systematic hardware-software co-design, the computationally intensive AlphaFold2 model can be efficiently accelerated on modern many-core CPUs equipped with dedicated matrix units. Our proposed framework, MatrixFold, significantly boosts inference performance by integrating specialized GEMM kernels, an optimized multi-NUMA communication strategy, an adaptive parallelism mechanism, and the first application of INT8 quantization for AlphaFold2 on a CPU platform.
Beyond AlphaFold2, the techniques proposed in MatrixFold are applicable to other Evoformer-based models such as RoseTTAFold and ESMFold, highlighting the broader impact of CPU-based quantized inference. 

Our evaluation reveals that MatrixFold is not only highly competitive in performance but also maintains prediction accuracy comparable to the FP32 baseline. Critically, for long protein sequences, its performance surpasses that of state-of-the-art GPU solutions on NVIDIA A800.
This study demonstrates that CPUs with matrix extensions are not only viable but can exceed top-tier GPUs in large-scale protein structure prediction workloads.
